\section{Working with GitHub}\label{sec:github}

GhostVault uses GitHub's REST API, version \code{2022-11-28}: the Git database endpoints (blobs,
trees, commits and references~\cite{gh-git}) for all of its data, the Contents API only to give an
empty repository its first commit and to look paths up one by one, and the repository endpoint to
check a repository and its permissions.

\subsection{Repositories and tokens}\label{sec:gh-checks}
A repository is added with its address (only \code{github.com} is accepted) and a token. GhostVault
accepts it only if GitHub describes it as private, not archived, and writable with that token
(\code{permissions.push}); a fine-grained token needs \emph{Contents: read and write} on it. When,
during an operation, GitHub rejects the token (401, 403) or no longer shows the repository, the
repository is marked as failing in Settings and new ghosts avoid it. A 404 for a branch is taken to
mean an empty repository only if the repository itself still answers, so a repository that has
become inaccessible is never mistaken for an empty one.

Tokens are kept in the operating system's keychain through the \code{keyring} library, which is
restricted to backends that store secrets in an operating-system keychain: Windows Credential
Manager, the macOS Keychain, Secret Service (libsecret) and KWallet. Plaintext or file-based backends
are never used, even when installed, and \code{keyring}'s fall-through chain is bypassed, so a
failing keychain can never fall through to a file. Without a usable keychain, tokens are kept in
memory for the session. Tokens never appear in GhostVault's files or log, and GitHub's own error
messages, where they are logged or kept in a record's notes, pass through a filter that removes
anything shaped like a token.

\subsection{Transport}
Requests go over one persistent HTTPS connection per thread to \code{api.github.com}, so a TCP and
TLS handshake is not paid per request. A proxy configured in the environment is honoured, with its
credentials, by tunnelling through it. Certificates are verified against the operating system's
trust store through \code{truststore}~\cite{truststore}, so certificate authorities installed by the
system or an organisation are trusted exactly as the system's browser trusts them; one TLS context is
passed explicitly to every request, and Python's process-wide defaults are never changed. Redirects
for renamed repositories are followed (301 and 302 for reads, 307 and 308 for any request), never to
another host and at most three in a row. A reply larger than 128\,MiB is not read.

Before an operation, a four-second check asks GitHub's \code{/rate\_limit} endpoint, authenticated
with a token, whether GitHub can be reached. An answer counts only if it comes from GitHub: a proxy
asking for a login, a sign-in page or filter answering in GitHub's place (no GitHub request id), and
a certificate that is not GitHub's are each reported as such instead of being taken for ``online''.

\subsection{Retries, rate limits and pacing}
Every request passes through one retry layer per vault (\tabref{tab:retries}). Rate limits are
waited out rather than treated as failures: a large upload can meet GitHub's limits, and waiting is
better than undoing hours of work. The window's progress line says until when it waits.

\begin{table}[t]
\caption{What the retry layer does}\label{tab:retries}
\centering\footnotesize
\begin{tabular}{@{}L{24mm}L{56mm}@{}}
\toprule
\textbf{Situation} & \textbf{What happens} \\
\midrule
no reply; 500, 502, 503, 504 & retried, if safe to repeat: at most 6 tries, exponential backoff with
full jitter (1, 2, 4\,s \dots, at most 60\,s) \\ \addlinespace[1.6pt]
429, or 403 with rate-limit signs & waited out: \code{Retry-After} (seconds or an HTTP date), else
until \code{X-RateLimit-Reset}, else 1, 2, 4 \dots 15 minutes; a single wait over 65 minutes, or
3 hours of waiting for one request, is reported instead \\ \addlinespace[1.6pt]
other 4xx & not retried: it describes a real problem \\ \addlinespace[1.6pt]
any write & serialised, at least 1\,s after the previous write \\
\bottomrule
\end{tabular}
\end{table}

GitHub asks integrators to make requests serially rather than concurrently, and to wait at least one
second between write requests~\cite{gh-best}. GhostVault serialises all writes and spaces them at
least one second apart. Downloads during a restore are reads and run up to eight at a time: below
GitHub's limit of 100 concurrent requests~\cite{gh-rate}, but not strictly serial; a rate-limit reply
to any of them is waited out like any other.

\subsection{Uncertain outcomes}\label{sec:uncertain}
Every read, and creating a blob, a tree or a commit, is safe to repeat: blobs are content-addressed,
and trees and commits are unreferenced until a branch points at them. Moving a branch and creating a
file through the Contents API are not. When one of those fails without a reply, or with an error
that does not say whether it took effect, the outcome is \emph{uncertain}, and GhostVault reads the
branch to find out: if it points at the new commit, the change happened; if it is unchanged, the
update is sent again (it can only fast-forward, so repeating it is safe); if someone else moved it,
the commit is rebuilt on top of the new head. A reply that names a different commit than the one
sent is resolved the same way. If GitHub refuses a tree because its base no longer holds a path being
deleted (a 422), the branch is read again and the commit rebuilt from it. A branch is never forced,
so another writer's commits to the same repository are never overwritten.

\subsection{Large repositories}
Listing a branch reads its tree recursively. GitHub truncates a recursive listing beyond 100\,000
entries or 7\,MB~\cite{gh-git}; GhostVault then looks each path it needs up on its own instead, which
is slower but correct. A repository of the size GitHub recommends, ``ideally less than 1\,GB, and
less than 5\,GB is strongly recommended''~\cite{gh-large}, holds at most about ten thousand chunks,
far below the point where listings are truncated.
